\subsection{Análisis de los indicadores financieros de rentabilidad: VAN, TIR, IR}

\textbf{Flujo de Efectivo}

Se elabora en base a los datos obtenidos en el presupuesto de caja (Figura 10.25). Primero colocamos el valor de la inversión que viene del plan de inversión (Figura 10.19). Para conocer el flujo neto de efectivo restamos el flujo de ingresos del flujo de egresos y en el año 6 sumamos la liquidación más cuentas x cobrar (Figura 10.26), liquidación inventarios finales (Figura 10.26) y el valor de salvamento (Figura 10.18) menos los pasivos circulantes (balance proyectado año 10 Tabla 35). Observar la (Figura 10.27)

\begin{table}[H]
\centering
\tiny
\renewcommand{\arraystretch}{1.15}
\setlength{\tabcolsep}{2pt}
\resizebox{\textwidth}{!}{%
\begin{tabular}{|p{3.2cm}|r|r|r|r|r|r|r|}
\hline
\multicolumn{8}{|c|}{\textbf{Elaboración Servicio}} \\
\multicolumn{8}{|c|}{\textbf{Flujo Neto de Efectivo}} \\
\multicolumn{8}{|c|}{\textbf{Original}} \\ \hline

\textbf{DESCRIPCIÓN} & \textbf{AÑO 0} & \textbf{AÑO 1} & \textbf{AÑO 2} & \textbf{AÑO 3} & \textbf{AÑO 4} & \textbf{AÑO 5} & \textbf{AÑO 6} \\ \hline

Inversión 
& L. 19,154.3 &  &  &  &  &  &  \\ \hline

Flujo de Ingresos 
& L. 0.0 & L. 140,163.5 & L. 282,416.4 & L. 486,816.1 & L. 710,504.4 & L. 1,015,908.9 & L. 1,437,879.4 \\ \hline

Flujo de Egresos 
&  & L. 211,144.4 & L. 276,106.2 & L. 340,374.2 & L. 632,135.5 & L. 655,886.2 & L. 774,795.8 \\ \hline

\textbf{flujo de efectivo} 
&  & \textbf{-S/ 70,980.87} & \textbf{S/ 6,310.24} & \textbf{S/ 146,441.88} & \textbf{S/ 78,368.84} & \textbf{S/ 360,022.75} & \textbf{S/ 663,083.62} \\ \hline

Más liquidación cuentas x cobrar 
&  &  &  &  &  &  & L. 116,535.84 \\ \hline

Más valor de salvamento 
&  &  &  &  &  &  & -14,625.00 \\ \hline

Menos pasivo circulante 
&  &  &  &  &  &  & -S/ 231,334.58 \\ \hline

\textbf{Flujo Neto de efectivo} 
& \textbf{S/ 19,154.29} & \textbf{-S/ 70,980.87} & \textbf{S/ 6,310.24} & \textbf{S/ 146,441.88} & \textbf{S/ 78,368.84} & \textbf{S/ 360,022.75} & \textbf{S/ 533,659.88} \\ \hline

\end{tabular}%
}

\vspace{0.2cm}
\textbf{Figura 10.27.} Flujo Neto de Efectivo \textbf{[Elaboración Propia]}
\end{table}

\textbf{Tasa de Corte}

Primero encontramos el costo de capital para ello necesitamos conocer cuanto representa el total de nuestra inversión respecto del financiamiento que utilizaremos 14.91\% y 85.09\% respectivamente y el valor de la tasa máxima a la cual podríamos invertir los fondos propios en el sistema bancario o financiero y la tasa mejor a la que lograríamos conseguir el financiamiento del préstamo, 6\% y 12\% respectivamente. Todo esto se resumen en la Figura 10.28
